% Caption for heatmap_fu_analytical_3panels.pdf  (Fig 8, fig:fu_heatmap)
% Panels: A IS, B pCS, C HA (same as five-panel / Fig 6).
% Y-axis is a multiplier of Deltombe CKD-5 fu_p, not fu_p itself.

\begin{figure}[htbp]
\centering
\includegraphics[width=\textwidth]{Figures/heatmap_fu_analytical_3panels.pdf}
\caption{%
\textbf{Analytical sensitivity of last-cycle PBUT plasma concentration to
unbound fractions.}
Panels \textbf{A} (indoxyl sulfate, IS), \textbf{B} ($p$-cresyl sulfate, pCS)
and \textbf{C} (hippuric acid, HA). Colour is the closed-form percentage
change $\Delta\overline{C}$ (\%) versus the no-drug CKD-5 baseline, not an
ODE map.
The \textbf{x}-axis is furosemide plasma unbound fraction $f_{u,p}$
(absolute units, $0.01$--$0.13$; the model value is $0.05$).
The \textbf{y}-axis is \emph{not} toxin $f_{u,p}$. It is a multiplier of the
CKD-5 unbound fraction reported by Deltombe et~al.\ (\textit{Toxins} 2015)
(\textbf{y}${}=1$):
IS $0.10$, pCS $0.05$, HA $0.52$.
Thus \textbf{y}${}=0.3$ means the toxin is more protein-bound than that
reference (IS $0.03$, pCS $0.015$, HA $0.16$),
and \textbf{y}${}=2$ means twice as unbound (IS $0.20$, pCS $0.10$, HA $1.04$).
The same y-scale is used on all three panels so one axis can cover toxins
with very different $f_{u,p}$.
The colour at $x=0.05$, $y=1$ is the analytical prediction at the paper's
usual binding values.
}
\label{fig:fu_heatmap}
\end{figure}
